\item Los astrónomos con frecuencia toman fotografías utilizando sólo la lente objetivo de un gran telescopio, sin emplear un ocular.
\begin{enumerate}
    \item Demuestre que el tamaño lineal $h'$ de la imagen formada por el objetivo está dado por $h' = f h / (f - p)$, donde $h$ es el tamaño del objeto, $f$ la distancia focal de la lente y $p$ la distancia al objeto.
    \item Simplifique esta expresión para el caso en el cual el objeto está a una distancia mucho mayor que la distancia focal del objetivo ($p \gg f$).
    \item La envergadura de la Estación Espacial Internacional (EEI) es de $108.6\text{ m}$. Determine el ancho de la imagen de la estación formada en el CCD por una lente objetivo de telescopio con $f = 4.00\text{ m}$ cuando orbita a una altitud de $407\text{ km}$.
\end{enumerate}